\section{Numerical validation and forecast}
\label{sec:numerics}

The theorem is analytic. Computation is used to attack the derivation, expose transitions, and convert a declared noise law into a shot budget; it is not used as a substitute for the proof.

\subsection{Adversarial theorem checks}

Independent scripts checked:
\begin{itemize}
  \item the one- and two-slot closed forms and transition equations through $n=12$;
  \item random product likelihoods against the flagged recovery inequality;
  \item random collective likelihood tables against \eqref{eq:collective};
  \item the exact coherent-accumulator circuit;
  \item asymmetric local strengths and fully adaptive outcome trees for $n=2,\ldots,5$.
\end{itemize}
The maximum numerical excess over the streaming theorem was $6.52\times10^{-14}$ for adaptive trees and $3.53\times10^{-14}$ for asymmetric products, consistent with numerical precision. The smallest tested collective-classical gap was positive. These checks do not prove the theorem or priority; they are regression tests against algebraic and implementation errors.

The release suite contains 25 automated tests. It verifies exact supports, transition behaviour, conservative systematic handling, fixed-sample planning, strict JSON output, deterministic simulation, and byte-stable same-environment figures and data.

\subsection{Transparent forecast model}

For planning only, we use the replaceable phenomenological law
\begin{equation}
\PA(n)=\frac{1+c_0c^n}{2},\qquad
\FR(n)=f_0 f^n.
\label{eq:forecast}
\end{equation}
The committed defaults are $c_0=f_0=0.995$, $c=0.998$, and $f=0.997$. They are illustrative logical-operation factors, not hardware calibration.

With $\lambda=1/2$, $\alpha=0.01$, target power $0.9$, $0.005$ systematic allowance in each branch, and $0.005$ null slack, the model's robust plan remains feasible through $n=143$ and becomes infeasible at $n=144$. The unadjusted expected score first crosses the $3/4$ asymptotic streaming null at $n=151$. Near the robust boundary the conservative Hoeffding shot count diverges rapidly, making uncertainty accounting more important than small circuit optimisations.

At $n=8$, a declared alternative $(\PA,\FR)=(0.97,0.95)$ under the same allowances needs 84 trials per branch for the Hoeffding guarantee ($\alpha=0.01$, power at least $0.9$). If 9700 of 10000 AUDIT trials and 9500 of 10000 RETURN trials succeed, the adjusted one-sided lower score is approximately $0.9443$, exceeding the enlarged $0.755$ null by approximately $0.1893$.

\begin{figure}[ht]
\centering
\includegraphics[width=\textwidth]{audit_return_benchmark.pdf}
\caption{Illustrative forecast from \eqref{eq:forecast}. Left: balanced score against the streaming, collective-classical, and robust planning thresholds. The one-slot streaming point is higher than the $n\geq2$ plateau. Right: conservative total trial count from \eqref{eq:hoeffding}; infeasible lengths are omitted. These curves are not hardware predictions.}
\label{fig:forecast}
\end{figure}
\FloatBarrier

\subsection{Functional coherent-memory length}

We define the \emph{certifiable functional coherent-memory length} as the largest preregistered $n$ for which \eqref{eq:certificate} is satisfied under a declared familywise testing rule. Unlike a bare $T_2$, it asks a memory to retain a temporal predicate and later participate in coherent global return. It remains benchmark-specific: changing the predicate, branch scores, or sequestration interface changes the resource being measured.
